\chapter{Инфраструктура проекта и тесты}
\label{ch:infra}
\index{Python}
\index{Pipenv}
\index{Docker}
\index{pytest}

\section{Окружение разработки}

Требуется \textbf{Python 3.12+}. Зависимости устанавливаются \textbf{только} в изолированное окружение (Pipenv по умолчанию).

\begin{itemize}
  \item \texttt{make install} --- создаёт \texttt{.venv/}, ставит пакеты из \texttt{Pipfile};
  \item команды в документации и Makefile --- через \texttt{pipenv run}.
\end{itemize}

\section{Структура репозитория}

\begin{tabularx}{\textwidth}{@{}l X@{}}
\toprule
\textbf{Каталог / файл} & \textbf{Назначение} \\
\midrule
\texttt{src\_starting\_point/} & Заготовка АБУ (отправная точка; намеренно неоптимальна как образец). \\
\texttt{src\_solution/} & Рабочая зона решения участника. \\
\texttt{external\_systems/digital\_mine/} & Прототип ЦР. \\
\texttt{external\_systems/regulator/} & Прототип Регулятора. \\
\texttt{tests/} & Интеграционные и прочие тесты репозитория. \\
\texttt{docs/} & Документация, диаграммы PlantUML. \\
\texttt{Makefile} & \texttt{install}, \texttt{tests-all}, пакет, сертификация, оценка, Docker. \\
\bottomrule
\end{tabularx}

\section{Тестирование}

Команда \texttt{make tests-all} запускает \texttt{pytest} и охватывает:

\begin{itemize}
  \item каталог \texttt{src\_starting\_point/tests/} --- юнит- и модульные тесты заготовки АБУ;
  \item корневой каталог \texttt{tests/} --- интеграционные и прочие тесты репозитория;
  \item конфигурацию \texttt{pytest.ini} в корне (в т.\,ч. маркеры);
  \item тесты безопасности с маркером \texttt{security} (отдельный прогон в Регуляторе с порогом покрытия критичного кода).
\end{itemize}

Прочие артефакты, связанные с проверками:

\begin{itemize}
  \item логи и отчёты \texttt{pytest} при сертификации в песочнице Регулятора;
  \item при необходимости --- \texttt{tests/test\_docker\_integration.py} и переменная \texttt{REGULATOR\_URL} для контейнерного сценария.
\end{itemize}

\section{Контейнеры}

Опционально:

\begin{itemize}
  \item \texttt{make docker-build}, \texttt{make docker-up};
  \item сервисы: порт 8080 (ЦР), 8081 (АБУ), 8082 (Регулятор);
  \item для интеграционных тестов с контейнером задаётся \texttt{REGULATOR\_URL}.
\end{itemize}

\section{Диаграммы PlantUML}

\begin{itemize}
  \item исходники: \texttt{docs/diagrams/*.puml};
  \item сборка PNG: \texttt{make diagrams} из корня репозитория (PlantUML, Java);
  \item вставка в PDF настоящего документа --- из \texttt{docs/diagrams/png/} (см. макрос \texttt{\textbackslash diagram} в преамбуле).
\end{itemize}
